\scanpage{36}
% archive sheet number: 35
\uline{Proposition 2}\quad Soit $f\in\operatorname{Aff}_B(G)$\deleted{\unsure{supposons}}\inserted{Supposons $G$ $A$-connexe}. Conditions équivalentes

\margin{$\unsure{X})$}

\[
\begin{aligned}
\text{(i)}\quad &f\in\operatorname{Aff}_A(G)\\
\text{(ii)}\quad &f\mathbin{\lrcorner} U_A(G)\subset U_A(G)\\
\text{(iii)}\quad &\langle f,U_A(G)\rangle\subset A
\end{aligned}
\]

Évidemment $\text{(i)}\Rightarrow\text{(ii)}\Rightarrow\text{(iii)}$ [\unsure{de|et} $\text{(}\deleted{\unsure{\mathrm{ii}}}\mathrm{iii}\text{)}\Rightarrow\text{(ii)}$ car \deleted{\unsure{soit}} $X\in U_A(G)$, pour prouver $f\mathbin{\lrcorner}X\in U_A(G)$ il suffit \inserted{(prop. 1)} de vérifier que $\langle g,f\mathbin{\lrcorner}X\rangle\in A$ pour $g\in\operatorname{Aff}_A(G)$, or $\langle g,f\mathbin{\lrcorner}X\rangle=\langle gf,X\rangle=\langle f,g\mathbin{\lrcorner}X\rangle\in\langle f,U_A(G)\rangle\subset A$] Prouvons $\text{(iii)}\Rightarrow\text{(i)}$ en utilisant que $G$ est $A$-connexe.

Soit $C\deleted{B}=\operatorname{Aff}_A(G)$, considérons le diagramme

\diagram{\begin{tikzcd}[ampersand replacement=\&]
C\arrow[r]\arrow[d,very thick] \& C_B\arrow[d]\\
\widehat C\simeq\operatorname{Hom}_A(U,A)\arrow[r] \& \widehat C_B=\operatorname{Hom}_B(U_B,B)
\end{tikzcd}}

où $\widehat{\phantom C}$ désigne le \inserted{\unsure{séparé}} complété pour l'idéal d'augmentation.
